% \subsubsection{Theorem and theorem-like environments}
% On the web, a theorem is emitted as a special |<div>|.
%    \begin{macrocode}
%<*classXimera>  

\ifdefined\ifxmtheoremenvtcolorbox
  % Already defined, presumably to set it to false
\else
  \expandafter\newif\csname ifxmtheoremenvtcolorbox\endcsname
  \xmtheoremenvtcolorboxtrue
\fi

\ifxmtheoremenvtcolorbox
    \usepackage{tcolorbox}
    \tcbuselibrary{breakable,skins}
    \tcbuselibrary{theorems}

    \tcbset{default-style/.style={
              boxrule     = 0pt,
              boxsep      = 0pt,
              enhanced jigsaw,
              sharp corners,
              before skip = 10pt,
              after skip  = 10pt,
              breakable,
      } 
    }
    
    \colorlet{theory}{green}
    \colorlet{example}{blue}
    \colorlet{remark}{yellow}
\fi


\ExplSyntaxOn

% ------------------------------------------------------------
% Logic for parsing of xmtheorem optional arguments
% ------------------------------------------------------------

\bool_new:N \l_xmtheorem_print_bool
\bool_new:N \l_xmtheorem_numbered_bool
\bool_new:N \l_xmtheorem_expandable_bool
\bool_new:N \l_xmtheorem_foldable_bool
\bool_new:N \l_xmtheorem_expanded_bool
\tl_new:N   \l_xmtheorem_title_tl
\tl_new:N   \l_xmtheorem_expandedstate_tl

\prop_new:N \g_xmtheorem_defaults_prop

\keys_define:nn { xmtheorem }
  {
    numbered .bool_set:N = \l_xmtheorem_numbered_bool,
    numbered .default:n = true,
    
    nonumbered .meta:n = { numbered=false },

    print .bool_set:N = \l_xmtheorem_print_bool,
    print .default:n = true,
    
    noprint .meta:n = { print=false },
    
    expandable .bool_set:N = \l_xmtheorem_expandable_bool,
    expandable .default:n = true,

    foldable .meta:n = { expandable=true, expanded=true },
    foldable .default:n = false,
    
    noexpandable .meta:n = { expandable=false },
    
    expanded .code:n =
      { \tl_set:Nn \l_xmtheorem_expandedstate_tl { open='' } },
    expanded .default:n = false,

    title .tl_set:N = \l_xmtheorem_title_tl,

    % Everything that isn't explicitly defined above is considered a title
    % (and the last wins...!)
    unknown .code:n = { \tl_set:NV \l_xmtheorem_title_tl \l_keys_key_tl }
  }

% ------------------------------------------------------------
% Global defaults
% ------------------------------------------------------------

\tl_new:N \g_xmtheorem_global_defaults_tl

\tl_gset:Nn \g_xmtheorem_global_defaults_tl
  {
    print=true,
    expandable=false,
    numbered=true
  }

% ------------------------------------------------------------
% Per environment defaults
% ------------------------------------------------------------

\NewDocumentCommand{\XmTheoremDefaults}{m m}
  {
    \prop_gput:Nnn \g_xmtheorem_defaults_prop {#1} {#2}
  }

  % Temporary variables/settings
\newif\ifxmcurrentexpandable
\newif\ifxmcurrentprint

\tl_new:N \l__xm_title_tl

\def\xmcurrenttitle{}  

% ------------------------------------------------------------
% Parser
% ------------------------------------------------------------

\cs_new_protected:Npn \xmtheorem_parse:nn #1 #2
  {
    % reset from global defaults
    \keys_set:nV { xmtheorem }
          \g_xmtheorem_global_defaults_tl
      
    % reset title
    \tl_clear:N \l_xmtheorem_title_tl

    % environment-specific defaults
    \prop_get:NnNT \g_xmtheorem_defaults_prop {#1} \l_tmpa_tl
      {
        \keys_set:nV { xmtheorem } \l_tmpa_tl
      }

    % local options override everything
    \clist_map_inline:nn {#2}
      {
            % We testen op $ en \ omdat titels vaak macro's of wiskunde bevatten.
            \regex_match:nnTF { \$ | \\ } {##1}
              {
                \tl_set:Nn \l_xmtheorem_title_tl {##1}
              }
              {
                % Geen wiskunde of macro's? Dan mag l3keys controleren of het een 
                % waarde-loze sleutel is (zoals 'numbered'), met je 'unknown' fallback.
                \keys_set:nn { xmtheorem } { ##1 }
              }
      }

      % set non-expl3 variables for actual use in LaTeX
      \bool_if:NTF \l_xmtheorem_print_bool
        { \xmcurrentprinttrue }
        { \xmcurrentprintfalse }

      \bool_if:NTF \l_xmtheorem_expandable_bool
        { \xmcurrentexpandabletrue }
        { \xmcurrentexpandablefalse }

      %% \xdef would cause tex to HANG when title contains eg \textit !!
      \tl_if_empty:NTF \l_xmtheorem_title_tl
          { \gdef\xmcurrenttitle{} } %  for \ifdefempty infra
          { \gdef\xmcurrenttitle{\tl_use:N \l_xmtheorem_title_tl} }
      \xdef\xmcurrentexpandedstate{\tl_use:N \l_xmtheorem_expandedstate_tl}      
  }

\NewDocumentCommand{\xmParseTheoremOptions}{m m}
{
   \xmtheorem_parse:nn {#1} {#2} 
  %  \bool_if:NTF \l_xmtheorem_print_bool {PRINT} {NOPRINT}
  %  \bool_if:NTF \l_xmtheorem_expandable_bool {EXPAND} {NOEXPAND}
  %  TITLE=\xmcurrenttitle.
}

\NewDocumentCommand{\initcap}{m}
  {
    \text_titlecase_first:n {#1}
  }

% ------------------------------------------------------------
% Styling: tcolorbox config
% ------------------------------------------------------------

\prop_new:N \g_xmtheorem_style_prop

% defaults

\tl_new:N \g_xmtheorem_style_default_base_tl
\tl_new:N \g_xmtheorem_style_default_backgroundpct_tl
\tl_new:N \g_xmtheorem_style_default_borderpct_tl

\tl_gset:Nn \g_xmtheorem_style_default_base_tl          { gray }
\tl_gset:Nn \g_xmtheorem_style_default_backgroundpct_tl { 5 }
\tl_gset:Nn \g_xmtheorem_style_default_borderpct_tl     { 45 }

% ------------------------------------------------------------
% Temporary variables used by the key parser
% ------------------------------------------------------------

\tl_new:N \l_xmtheorem_style_base_tl
\tl_new:N \l_xmtheorem_style_backgroundpct_tl
\tl_new:N \l_xmtheorem_style_borderpct_tl

% ------------------------------------------------------------
% Keys
% ------------------------------------------------------------

\keys_define:nn { xmtheoremenv/config }
  {
    color_base          .tl_set:N = \l_xmtheorem_style_base_tl,
    color_backgroundpct .tl_set:N = \l_xmtheorem_style_backgroundpct_tl,
    color_borderpct     .tl_set:N = \l_xmtheorem_style_borderpct_tl,
  }

% ------------------------------------------------------------
% User commands
% ------------------------------------------------------------

\NewDocumentCommand{\xmSetTheoremStyle}{mm}
  {
    % start from defaults
    \tl_set_eq:NN
      \l_xmtheorem_style_base_tl
      \g_xmtheorem_style_default_base_tl

    \tl_set_eq:NN
      \l_xmtheorem_style_backgroundpct_tl
      \g_xmtheorem_style_default_backgroundpct_tl

    \tl_set_eq:NN
      \l_xmtheorem_style_borderpct_tl
      \g_xmtheorem_style_default_borderpct_tl

    % apply user keys
    \keys_set:nn { xmtheoremenv/config } { #2 }

    % store
    \prop_gput:NnV
      \g_xmtheorem_style_prop
      { #1/base }
      \l_xmtheorem_style_base_tl

    \prop_gput:NnV
      \g_xmtheorem_style_prop
      { #1/backgroundpct }
      \l_xmtheorem_style_backgroundpct_tl

    \prop_gput:NnV
      \g_xmtheorem_style_prop
      { #1/borderpct }
      \l_xmtheorem_style_borderpct_tl
  }

\NewDocumentCommand{\xmSetDefaultTheoremStyle}{m}
  {
    \keys_set:nn { xmtheoremenv/config } { #1 }

    \tl_gset_eq:NN
      \g_xmtheorem_style_default_base_tl
      \l_xmtheorem_style_base_tl

    \tl_gset_eq:NN
      \g_xmtheorem_style_default_backgroundpct_tl
      \l_xmtheorem_style_backgroundpct_tl

    \tl_gset_eq:NN
      \g_xmtheorem_style_default_borderpct_tl
      \l_xmtheorem_style_borderpct_tl
  }

% ------------------------------------------------------------
% Accessors with fallback to defaults
% ------------------------------------------------------------

\cs_new:Npn \xmtheorem_style_base:n #1
  {
    \prop_if_in:NnTF \g_xmtheorem_style_prop { #1/base }
      { \prop_item:Nn \g_xmtheorem_style_prop { #1/base } }
      { \tl_use:N \g_xmtheorem_style_default_base_tl }
  }

\cs_new:Npn \xmtheorem_style_backgroundpct:n #1
  {
    \prop_if_in:NnTF \g_xmtheorem_style_prop { #1/backgroundpct }
      { \prop_item:Nn \g_xmtheorem_style_prop { #1/backgroundpct } }
      { \tl_use:N \g_xmtheorem_style_default_backgroundpct_tl }
  }

\cs_new:Npn \xmtheorem_style_borderpct:n #1
  {
    \prop_if_in:NnTF \g_xmtheorem_style_prop { #1/borderpct }
      { \prop_item:Nn \g_xmtheorem_style_prop { #1/borderpct } }
      { \tl_use:N \g_xmtheorem_style_default_borderpct_tl }
  }

% ------------------------------------------------------------
% xcolor specifications
% ------------------------------------------------------------

\cs_new:Npn \xmtheorem_style_backspec:n #1
  {
    \xmtheorem_style_base:n {#1}
    !
    \xmtheorem_style_backgroundpct:n {#1}
    !
    white
  }

\cs_new:Npn \xmtheorem_style_borderspec:n #1
  {
    \xmtheorem_style_base:n {#1}
    !
    \xmtheorem_style_borderpct:n {#1}
    !
    black
  }

% Temp variable 
\str_new:N \g_xm_postcode_name_str

% Create a (pdf-version of) a 'newtheorem' with extra print=false or expandable=true optional arguments
%  This is probably only to be used by the wrapper \xmNewTheoremEnv  (that also creates an HTML version !)
%  TODO(?): undef a possible previous version, so that it can also re-define an environment ? 
\NewDocumentCommand{\xmNewTheoremEnvPDF}{m m}
  {
    % Create a (presumably) 'standard' amsthm-theorem; THIS USES THE CURRENTLY ACTIVE (amsthm-)STYLE
    \newtheorem{#1}{#2}
    
    %\newcommand{\theH#1}{\thetitlenumber.\the#1}   % TODO: check if needed/wanted for hyperref ...
    
    % Allow for numbering 1,2 vs 2.1.1, 2.1.2 to depend on in-xourse-or-not.
    \cs_set:cpn { the#1 }
    {
      \xmtheoremtitle{#1}
    }

    % Save the original macro's, to be able to create a wrapper around the amsthm-version with the original name, but with extra optional arguments
    \cs_gset_eq:cc { #1@inner }    { #1 }
    \cs_gset_eq:cc { end#1@inner } { end#1 }

    \cs_if_exist:cT { c@#1 }
      {
        \newcounter{#1@inner}
      }
    
    % Setup styles globally; only the color can be changed for now (but with \tcbset everything is possible )
    \ifxmtheoremenvtcolorbox   
        \tcbset{#1-style/.style={
                  default-style,
                  colback = \xmtheorem_style_backspec:n{#1} ,
                  borderline~west = { 2pt }{ 0pt }{ \xmtheorem_style_borderspec:n{#1}} ,
             } 
        }
    \fi

    % Replace original with wrapper
    \RenewDocumentEnvironment{#1}{ +O{} }
      {
        \xmtheorem_parse:nn {#1} {##1}%

        \def\xmthmprecode{precode#1}
        \def\xmthmpostcode{postcode#1}


        \str_gset:Nn \g_xm_postcode_name_str { postcode #1 }    % save for later use
        
        %  \typeout{DBG:~reneweddocenv~#1~with~\l_xmtheorem_title_tl.}

        \bool_if:NF \l_xmtheorem_numbered_bool
          {
            % No numbers, and 'eat the blank' with \!.
            \cs_set:cpn { the#1 } {\!}
            % When no number, also not counted; so undo the add that already happened
            \addtocounter{#1}{-1}
          }        

        \bool_if:NTF \l_xmtheorem_print_bool
          {
            \ifxmtheoremenvtcolorbox
              % Add the tcolorbox env (with a macro because of teX-nesting-issues)
              \csname tcolorbox\endcsname[#1-style]%
            \fi

            % Add the (original) amsthm-env (with a macro because of teX-nesting-issues)
            \tl_if_empty:NTF \l_xmtheorem_title_tl
                {\csname #1 @inner \endcsname }
                {\csname #1 @inner \endcsname [ \l_xmtheorem_title_tl ]}
          }
          {%
            % Do not print (unless xmDebug is defined)
            \cs_if_exist:cTF { xmDebug } 
              {             \bgroup\color{red} }%
              {\setbox0\vbox\bgroup} % HACK: print to hidden box0 !!!%   
                \csname #1 @inner \endcsname%
                %%% WRONG \setcounter{#1}{\value{#1@inner}}%
          }
        % Execute \precode<envname>  if it is defined
        %%\cs_if_exist:cT { precode #1 } {\csname precode #1 \endcsname}
        \expandafter\ifcsname \xmthmprecode \endcsname \csname\xmthmprecode\endcsname \fi        
      }
      {% At END of environment
        % Execute \postcode<envname>  if it is defined
        \expandafter\ifcsname \xmthmpostcode \endcsname \csname\xmthmpostcode\endcsname \fi        
        %%\cs_if_exist_use:c { \g_xm_postcode_name_str }
        %%%  This DOES NOT WORK, thus used temp variable !!! 
        %%%     \cs_if_exist:cT { postcode #1 } {\csname postcode #1 \endcsname}
        \bool_if:NTF \l_xmtheorem_print_bool
          {
            % Close the tcolorbox environment if it was opened
            \ifxmtheoremenvtcolorbox
              \csname endtcolorbox\endcsname%
            \fi
          }
          {%
            % Close the non-print extra bgroup
            \egroup%
          }
      }
    % \typeout{DBG:~Renewed~environment~#1}
  }


%% To use infra in explsyntaxoff :-(
\NewDocumentCommand{\xmTheoremNoNumber}{m}
  {
        \bool_if:NF \l_xmtheorem_numbered_bool
          {
            \cs_set:cpn { the#1 } {}
            \addtocounter{#1}{-1}
          }
  }

\ExplSyntaxOff

\providecommand{\xmTheoremHeadLevel}{%
  \romannumeral-`\q% Dit is een truc die TeX dwingt alles te expanderen tot een getal en de 'q' wist
  \ifnum\value{subsubsection}>0   6%
  \else\ifnum\value{subsection}>0 5%
  \else                           4%
  \fi%
  \fi%
}

\newcommand{\xmNewTheoremEnvHTML}[2][]{%

  \newcounter{#2}
  \newenvironment{#2}[1][]%
  {%
    % Parse optional argument locally
    \xmParseTheoremOptions{#2}{##1}    %% environment and optional args
%
    \def\xmthmprecode{precode#2}
    \def\xmthmpostcode{postcode#2}
%
    \refstepcounter{problem}%
    \stepcounter{identification}%
%
    % update theorem counter if it exists
    \ifcsdef{c@#2}
      {\refstepcounter{#2}}
      {}%
%
    \xmTheoremNoNumber{#2}  % Removes the number is nonumbered
%
    \ifvmode\IgnorePar\fi
    \EndP
%
XXX
% 
      \ifxmcurrentexpandable
        \HCode{
          <section>
          <details
          class="theorem-like problem-environment #2"
          id="problem\arabic{identification}"
          \xmcurrentexpandedstate >
        }
        %% DBG: STATE=\xmcurrentexpandedstate   % \HCode might silently NOT print it !!!
%
        \HCode{<summary class="summary-#2">}%
%
        \HCode{<h\xmTheoremHeadLevel }
        \HCode{ class="summary-header">}%
          \initcap{\GetTranslation{#1}}
          \ifcsvoid{the#2}%
            {}
            {\csuse{the#2}.}%
        \HCode{</h\xmTheoremHeadLevel>}
        %
        \ifdefempty{\xmcurrenttitle}
          {}
          {
            \HCode{<span class="summary-title">}%
            \xmcurrenttitle%
            \HCode{</span>}%
          }
        \HCode{</summary>}
      \else
        \HCode{
          <section>
          <div
          class="theorem-like problem-environment #2"
          id="problem\arabic{identification}"
          titletext="\GetTranslation{#1}">
        }
        \HCode{<h\xmTheoremHeadLevel }
              \HCode{ class="summary-header">}%
          \initcap{\GetTranslation{#1}}
          \ifcsvoid{the#2}%
            {}
            {\csuse{the#2}.}%
        \HCode{</h\xmTheoremHeadLevel>}
%
        \ifdefempty{\xmcurrenttitle}
          {}
          {
            \HCode{<span class="theorem-like-title">}%
            \xmcurrenttitle%
            \HCode{</span>}%
          }
      \fi
    \IgnoreIndent
    \expandafter\ifcsname \xmthmprecode \endcsname \csname\xmthmprecode\endcsname \fi
  }
  {
    \expandafter\ifcsname \xmthmpostcode \endcsname \csname\xmthmpostcode\endcsname \fi
    \ifvmode\IgnorePar\fi
    \EndP
      \ifxmcurrentexpandable
        \HCode{</details></section>}
      \else
        \HCode{</div></section>}
      \fi
    \IgnoreIndent
  }%
}

\newcommand{\xmNewTheoremEnv}[2][]{%
  \ifdefined\HCode
    \xmNewTheoremEnvHTML[#1]{#2}
  \else
    \xmNewTheoremEnvPDF{#2}{\initcap{\GetTranslation{#1}}}
  \fi
}


% if in xourse, then prepend part/chapter/etc (as found in \titlenumber) 
\newcommand{\xmtheoremtitle}[1]{%
    \ifdefined\isXourse\thetitlenumber.\arabic{#1}\else\arabic{#1}\fi%
}

% OBSOLETE, use \xmtheoremtitle
\providecommand{\theoremtitle}[1]{%   %% used in printstyle ...
    \ifdefined\isXourse\thetitlenumber.\arabic{#1}\else\arabic{#1}\fi%
}


% ------------------------------------------------------------
% Default configurations for newtheorems
% ------------------------------------------------------------


\xmSetDefaultTheoremStyle{ color_base=gray,    color_backgroundpct=5,  color_borderpct=45 }

\xmSetTheoremStyle{definition}  { color_base=theory,  color_backgroundpct= 5, color_borderpct=45 }
\xmSetTheoremStyle{notation}    { color_base=theory,  color_backgroundpct= 5, color_borderpct=45 }
\xmSetTheoremStyle{corollary}   { color_base=theory,  color_backgroundpct= 5, color_borderpct=45 }
\xmSetTheoremStyle{axiom}       { color_base=theory,  color_backgroundpct= 5, color_borderpct=45 }
\xmSetTheoremStyle{theorem}     { color_base=theory,  color_backgroundpct= 2, color_borderpct=65 }
\xmSetTheoremStyle{proposition} { color_base=theory,  color_backgroundpct= 2, color_borderpct=65 }
\xmSetTheoremStyle{corollary}   { color_base=theory,  color_backgroundpct= 2, color_borderpct=65 }
\xmSetTheoremStyle{algorithm}   { color_base=theory,  color_backgroundpct= 2, color_borderpct=65 }
\xmSetTheoremStyle{explanation} { color_base=theory,  color_backgroundpct= 1, color_borderpct=75 }

\xmSetTheoremStyle{example}     { color_base=example, color_backgroundpct= 3, color_borderpct=75 }
\xmSetTheoremStyle{exercise}    { color_base=example, color_backgroundpct= 0, color_borderpct=50 }

\xmSetTheoremStyle{remark}      { color_base=remark,  color_backgroundpct= 5, color_borderpct=75 }
\xmSetTheoremStyle{warning}     { color_base=remark,  color_backgroundpct= 2, color_borderpct=65 }

% Set default amsthm style
\theoremstyle{definition} % No italic (because this makes also text in TikZ italic !!!)


%% Create all Ximera-provided theorem-environments
\ExplSyntaxOn
\seq_set_from_clist:Nn \l_my_theorems_seq
  { 
theorem,
algorithm,
axiom,
claim,
conclusion,
condition,
conjecture,
corollary,
criterion,
definition,
example,
explanation,
fact,
lemma,
formula,
idea,
notation,
model,
observation,
proposition,
paradox,
procedure,
remark,
summary,
template,
warning
  }

\seq_map_inline:Nn \l_my_theorems_seq
  {
    \xmNewTheoremEnv[#1]{#1}    % TODO: fix optional title ...
  }

\ExplSyntaxOff

% % TODO: implement  non-numbered environments (cfr \newtheorem* ... )!
\ifdefined\HCode
   \let\c@explanation@inner\relax
\else
  \renewcommand{\theexplanation}{}   % HACK: no numbers shown
\fi

%</classXimera>
% 
%    \end{macrocode}
%    \begin{macrocode}
%<*htXimera>
% % OBSOLETE (?) See \xmNewTheoremEnv
% 
\newcommand{\ConfigureTheoremEnv}[1]{%
\renewenvironment{#1}[1][]{\refstepcounter{problem}%
\ifthenelse{\equal{##1}{}}{}{%
  \HCode{<span class="theorem-like-title">}##1\HCode{</span>}%
}}{}
\ConfigureEnv{#1}{\stepcounter{identification}\ifvmode \IgnorePar\fi \EndP\HCode{<div class="theorem-like problem-environment #1" id="problem\arabic{identification}" titletext=" \GetTranslation{#1}">}}{\ifvmode\IgnorePar\fi\EndP\HCode{</div>}\IgnoreIndent}{}{}%
}
%</htXimera>
%    \end{macrocode}
